\section{Discussion}

In the bank's microlearning program, response accuracy increased with each return of a question, also
in the cohort in which the same sequences were compared. The increase resulted from error corrections
outnumbering forgetting events. Some errors, however, persisted despite feedback and consisted in
choosing the same distractor again.
The direction of change is consistent with studies of spaced education in medical training.
In urology residents, the knowledge advantage gained from periodically recurring questions persisted
two years after the end of the program \citep{kerfoot2009}, and a systematic review of studies on
clinicians' continuing professional development indicated a beneficial effect of this method on the
acquisition and retention of knowledge \citep{phillips2019}.

The conditions of those studies, however, differed from those of the present analysis. The
participants were health professionals, and the effect was usually assessed with a test administered
outside the program, in comparison with a control group. In the program studied here, the participants
were bank employees, and the increase in accuracy was observed only in responses given within the
program itself. The agreement therefore concerns the direction of change, not the size of the effect or
its cause. The results do, however, show what the increase in accuracy consisted of, which a final test
score does not reveal. It was the difference between two opposing processes, error correction and
forgetting (\NCorrections{} vs.\ \NForgetting{} events), and the same overall increase could also
arise with more frequent forgetting if it were accompanied by a correspondingly larger number of
corrections. Aggregate measures also do not detect errors that persist despite feedback.

Forgetting was associated with three factors. The first was the passage of time: forgetting was more
frequent at longer intervals, in line with research on the spacing between repetitions
\citep{cepeda2006}. The second was a short history of correct responses. Knowledge confirmed by one
correct response was forgotten two to four times more often than knowledge confirmed several times in
a row, at every interval.
Similarly, in studies of repeated practice, each additional successful retrieval increased the
durability of knowledge \citep{rawson2013}. The third factor was error relapse. Error correction proved
the least durable, and when it did not hold, the employee usually returned to the distractor chosen
before the correction. The return of original errors after feedback has also been described in
laboratory studies \citep{butler2011}. The correct answer may thus not have replaced the misconception
but coexisted with it.

Repeated errors differed in character from forgetting. Employees chose the same distractor more often
than chance would predict, regardless of the interval, and each successive error reduced the chance of
correction. Such a pattern is consistent with a persistent misconception rather than with a lack of
knowledge, which would lead to guessing. A selected distractor can become fixed in memory, and feedback
reduces this effect \citep{butler2008}. In the program studied, the feedback gave the correct answer
and an explanation but did not address the selected distractor, although explaining why the selected
answer is wrong is usually more effective \citep{shute2008}. The concentration of repeated errors in
\PersistTopItems{} questions further suggests that some of them resulted from question design, e.g.,
an ambiguous distractor or content inconsistent with the bank's practice. Such questions can be
identified by standard distractor analysis \citep{gierl2017}.

In one-off training, e.g., in information security, the effects faded after several months
\citep{reinheimer2020}. The present results indicate that knowledge requires further repetition also in
a program with periodically recurring questions, and that the way these repetitions are scheduled can
be improved. Because the interval did not depend on the outcome, mastered questions returned as often
as questions in which the employee had only just corrected an error. An adaptive schedule could
lengthen the interval after several consecutive correct responses and shorten it after an error
correction, when the risk of relapse is highest. Personalized review improved long-term retention
compared with a fixed schedule \citep{lindsey2014}, and models of forgetting rate based on response
history are already used in language learning applications \citep{settles2016}. A second change could
be feedback addressing the selected distractor, introduced first in the questions with the largest
number of repeated errors.

\subsection{Limitations and future research}
The study was observational and had no control group, so the increase in accuracy cannot be attributed
solely to the program, as employees also acquired knowledge in their everyday work. Because the
questions were repeated, the data do not allow consolidation of knowledge to be distinguished from
memorization of the answer to a specific question. Response time and employee confidence were not
recorded, so forgetting cannot be separated from inattention or from a lucky guess at an earlier
attempt, which learner models capture as slip and guess \citep{corbett1994}. Such a source of part of
the forgetting is suggested by its rate of \ForgetShort\% already at intervals shorter than 14 days.
The longest intervals may have included periods when the program was suspended, so changes in question
content may have inflated the forgetting rate at these intervals. The comparisons were descriptive and
were not corrected for multiple testing. The confidence intervals accounted for the dependence between
responses of the same employee, but not of the same question. The associations with the number of
consecutive correct responses or errors may therefore partly reflect differences in difficulty between
questions.

Future studies could compare, in a randomized design, a fixed with an adaptive schedule and standard
feedback with feedback addressing the selected distractor, using equivalent versions of questions that
would separate knowledge from memorization of the answer. Recording response confidence would make it
possible to test whether persistent errors concern high-confidence beliefs \citep{butler2011}.
Microlearning outcomes should also be linked to the application of knowledge at work, e.g., to the
number of procedural errors \citep{blume2010}.

\section{Conclusions}
In a program with periodically recurring questions, the response accuracy of bank employees increased
from attempt to attempt. Knowledge nevertheless required further repetition, because it was forgotten
over time, particularly after the first correct response and after an error correction. Some errors
persisted despite feedback and consisted in choosing the same distractor again. A schedule adapted to
employee performance and feedback addressing the selected distractor may increase the effectiveness of
such programs, but their effect requires confirmation in studies with a control group.
